\section{Nichtfunktionale Anforderungen}
Neben den funktionalen und technischen Anforderungen spielen nichtfunktionale Eigenschaften eine wesentliche Rolle für die Qualität des Systems.
Die Zuverlässigkeit des Systems ist insbesondere im Dauerbetrieb von zentraler Bedeutung.
Gleichzeitig wird eine modulare Software Architektur angestrebt, um den Code wiederverwendbar und strukturiert aufzubauen, was die Wartbarkeit sicherstellen soll.
Alle nichtfunktionale Anforderungen sind in Tabelle~\ref{tab:Nichtfunktionale Anforderung} aufgelistet.

\begin{table}[H]
\centering
\begin{tabular}{p{3cm} p{9cm}}
\textbf{ID} & \textbf{Beschreibung} \\
\hline
N1 & Hohe Zuverlässigkeit im Dauerbetrieb \\
N2 & Energieeffizienter Betrieb trotz hoher LED-Anzahl \\
N3 & Wartbarkeit und klare Struktur des Codes \\
N4 & Ansprechendes und integrierbares Design \\
N5 & Benutzerfreundliche Bedienbarkeit \\
\end{tabular}
\caption{Nichtfunktionale Anforderungen}
\label{tab:Nichtfunktionale Anforderung}
\end{table}



\subsection{Optische Anforderungen}
Gestalterische Aspekte spielen ebenfalls eine Rolle, da die Wortuhr sowohl als technisches System als auch als Gestaltungselement eingesetzt wird.
Der wichtigste Aspekt ist die Lesbarkeit und gleichmäßige Ausleuchtung der Anzeige. 
Die Worte müssen aus typischen Betrachtungsabständen klar erkennbar und ausreichend stark beleuchtet sein, dass ein Kontrast zwischen aktiven und inaktiven Buchstaben entsteht.
Weiterhin ist eine homogene Lichtverteilung sicherzustellen, um Hotspots oder blenden beim betrachten zu vermeiden.
Die Hauptanzeige soll in einem Raster mit gleicher Seitenlänge angeordnet sein.
Alle optischen Anforderungen sind in Tabelle~\ref{tab:Optische Anforderung} aufgelistet.

\begin{table}[H]
\centering
\begin{tabular}{p{3cm} p{9cm}}
\textbf{ID} & \textbf{Beschreibung} \\
\hline
O1 & Gewährleistung einer guten Lesbarkeit der Anzeige bei typischen Betrachtungsabständen \\
O2 & Ausreichende Helligkeit der aktiven Buchstaben zur Erzeugung eines klaren Kontrasts zu inaktiven Elementen \\
O3 & Gleichmäßige Ausleuchtung der gesamten Anzeige ohne lokale Überhellungen \\
O4 & Vermeidung von Hotspots und Blendwirkungen bei der Betrachtung \\
O5 & Homogene Lichtverteilung über alle dargestellten Wörter \\
O6 & Anordnung der Hauptanzeige in einem gleichmäßigen Raster mit identischer Seitenlänge \\
\end{tabular}
\caption{Funktionale Anforderungen der Wortuhr}
\label{tab:Optische Anforderung}
\end{table}

% %%%%%%%%%%%Old Version
% Die nicht-funktionalen Anforderungen beschreiben qualitative Eigenschaften des Systems.
% Das System soll eine hohe Zuverlässigkeit aufweisen. Die mittlere Betriebsdauer zwischen zwei Ausfällen soll mindestens 1000~h betragen.
% Die Wartbarkeit soll durch eine strukturierte Softwarearchitektur sowie durch dokumentierte Schnittstellen unterstützt werden. Der Austausch einzelner Hardwarekomponenten soll ohne vollständige Demontage möglich sein.
% Die Konstruktion soll modular aufgebaut sein. Anzeigeeinheit, Steuerplatine und Sensorik sollen getrennt ausgeführt werden. Erweiterungen, beispielsweise zusätzliche Informationsanzeigen, sollen ohne grundlegende Systemänderung integrierbar sein.
% Die elektrische Sicherheit soll den geltenden Normen für Niederspannungsgeräte entsprechen. Die Erwärmung der Bauteile soll innerhalb zulässiger Grenzwerte bleiben.
